\begin{lemma}[Situation 7.1, Claim 9]
Fix a minimum hierarchy-counterexample context and suppose
\[
 \noderef{HierarchyPathNumber}(C.\mathit{hierarchy})=1.
\]
Assume explicitly the paper's ambient palette-surplus consequence: for
every proper replacement coloring leaving exactly \(e_0\) uncolored,
every graph vertex has a missing color.  Then the finite set
\noderef{HierarchyEligibleCommonMissingIndices} is nonempty.  Equivalently,
for some \(1\leq i\leq\ell-1\), there are a full replacement presentation
in the paper's same-rung minimum-counterexample family \(\mathcal K\) and
a color missing both at the exact tail vertex \(y_i\) and at the final
vertex \(y_\ell\).
\end{lemma}
\begin{proof}
We first exhibit the original presentation as an actual member of the
replacement family.  Let \(\mathcal S\) be the certified series of the
given ETT.  In the structure
\noderef{CorrespondingExtendedTashkinovTree}, set the replacement input
equal to the original input, the replacement series equal to
\(\mathcal S\), and the replacement ETT equal to the given ETT.  The
rung, dependent-series, ladder-tree, stored-set, connecting-edge, and
extension-kind agreements are then literal reflexivities.  The required
final-color equation is not a reflexivity: it is exactly the
\(\mathit{generatingColoring}\) field of
\noderef{ExtendedTashkinovTree}.

We verify the remaining coloring fields at every series index.  By
\noderef{TashkinovSeries}, each index has a certified
\noderef{TashkinovSeriesStep}.  At index zero its initial constructor
identifies the indexed coloring with the input coloring, and its
\noderef{KTriple} field says that this input is proper and leaves exactly
\(e_0\) uncolored.  At a positive SE index the constructor identifies the
current coloring with the selected coloring in its
\noderef{MaximumDefectiveChoice}; the
\(\mathit{selectedProper}\) and \(\mathit{selectedOneUncolored}\) fields
give the two required assertions.  At a PE index they are the
\(\mathit{currentProper}\) and \(\mathit{currentOneUncolored}\) fields of
the constructor.  At an RE index the \(\mathit{currentColoring}\) equation
identifies the current coloring with the preceding one, to which the
\(\mathit{previousProper}\) and
\(\mathit{previousOneUncolored}\) fields apply.  These four constructors
exhaust all certified steps.

At every positive series index the replacement and original colorings are
the same function.  Both clauses in \noderef{StableColoring} are therefore
checked directly: every controlled incident edge keeps its value by
reflexivity, and at every vertex of the ladder tree the two missing-color
sets are equal by reflexivity.  This finishes the identity
correspondence certificate.  Reuse the \(\mathit{maximum}\) field of
\noderef{HierarchyCounterexample}, the original base and counterexample,
and the supplied \noderef{HierarchyMinimumCounterexampleContext}.
Finally,
\noderef{SameGoodHierarchy} holds reflexively: its underlying
\noderef{SameHierarchyStages} data, every transported stored set, and
every relevant Gamma set are unchanged.  The fields just listed form a
\noderef{HierarchyReplacementPresentation} \(P_0\) whose coloring is
\(\varphi_n\) and whose terminal tree is the same exact edge-valued
tree \(T\).

Write
\[
 H_0=C.\mathit{hierarchy},\qquad
 q=H_0.\mathit{dividerCount},\qquad
 A=\noderef{HierarchyStarredStage}(H_0,\operatorname{Fin.last}q),
 \qquad
 \ell=\noderef{HierarchyFullTailLength}(C).
\]
Apply \noderef{HierarchyCounterexampleTerminalCommonMissing} to the
properness, one-uncoloredness, lower-rung package, fewer-divider
induction, lexicographic minimality, and large-missing-set fields of the
supplied \noderef{HierarchyMinimumCounterexampleContext}.  It gives a
color \(\alpha\) such that
\[
 \alpha\in\overline{\varphi_n}(B),\qquad
 \alpha\in\overline{\varphi_n}(y_\ell),\qquad
 B=\noderef{HierarchyFullTailPrefix}(C,\ell-1).               \tag{1}
\]
The same fields, through
\noderef{HierarchyCounterexampleTailAtLeastTwo}, give
\[
                              2\leq\ell.                      \tag{2}
\]

If some \(i\) with \(1\leq i\leq\ell-1\) and some color \(\xi\)
are missing both at the last vertex of
\(\noderef{HierarchyFullTailPrefix}(C,i)\) and at \(y_\ell\),
then \(i,P_0,\xi\) are exactly the witnesses in the filter defining
\noderef{HierarchyEligibleCommonMissingIndices}; this proves the
conclusion.  We may consequently assume that, for every such \(i\), the
two missing-color sets are disjoint.                              \(\tag{3}\)

We now locate the first witness in (1) exactly.  First,
\(A\) is an ordered initial segment of \(T\).  If \(q=0\), the
zero-index branch of \noderef{HierarchyStarredStage} makes \(A\) the
base core, and the \(\mathit{coreInsideFirst}\) and
\(\mathit{terminalStage}\) fields of
\noderef{TashkinovHierarchy} give \(A\preccurlyeq T\).  If \(q>0\), the
positive-index branch makes \(A\) the last ordinary hierarchy stage, and
the last \(\mathit{nested}\) relation together with
\(\mathit{terminalStage}\) gives the same conclusion.  These are exact
ordered initial segments in the sense of
\noderef{TreeSequenceInitialSegment}, so every vertex and individual edge
entry agrees at its numerical index.

By \noderef{HierarchyFullTailLength}, \(\ell=|T|-|A|\).
The length inequality in \(A\preccurlyeq T\) removes the truncated
subtraction and gives
\[
                         |A|+\ell=|T|.                        \tag{4}
\]
Using (2) and (4), the minimum in
\noderef{HierarchyFullTailPrefix} evaluates to
\[
 |B|=\min\{|A|+\ell-1,|T|\}=|T|-1=|A|+\ell-1.                \tag{5}
\]
The same definition identifies every entry of \(B\) with the entry of
\(T\) at the same index, and the common-extension clauses of
\noderef{TreeSequenceInitialSegment} give \(A\preccurlyeq B\).

Expand the first membership in (1) using
\noderef{TreeMissingColors}.  It supplies a vertex \(v\) of \(B\) at
which \(\alpha\) is missing.  By
\noderef{TreeSequenceVertexSet}, \(v\) is the vertex entry of \(B\) at
some index \(r\) with \(0\leq r\leq |B|\).  If \(r\leq |A|\), the exact
initial-segment equality \(A\preccurlyeq B\) puts \(v\) in \(V(A)\).
If \(r>|A|\), set \(i=r-|A|\).  Equations (2) and (5) give
\(1\leq i\leq\ell-1\), and the entrywise prefix equation says that
\(v\) is precisely the displayed vertex \(y_i\), namely the vertex of
\(T\) at position \(|A|+i\).  But \(\alpha\) is also missing at
\(y_\ell\) by (1), contradicting (3).  Therefore the first case holds,
and
\[
             \alpha\in\overline{\varphi_n}(A),\qquad
             \alpha\in\overline{\varphi_n}(y_\ell).           \tag{6}
\]

We next obtain the manuscript's normalization.  Let \(n\) be the last
series index of the given ETT.  Its \(\mathit{roundsPositive}\) field in
\noderef{ExtendedTashkinovTree} says that \(n\) is a positive index.
Inspecting the certificate at \(n\) from
\noderef{TashkinovSeries}, the initial constructor of
\noderef{TashkinovSeriesStep} is therefore impossible; its
\(\mathit{currentKind}\) field in the remaining constructors says that
the terminal kind is exactly one of SE, PE, and RE.

If \(q>0\), the positive branch of
\noderef{HierarchyStarredStage} identifies \(A\) with \(T_{n,q}\), so
(6) is the required positive-divider normalization.  Suppose \(q=0\).
For terminal kind SE or RE, the nonparallel clause of
\noderef{HierarchyBase} identifies the base core, hence \(A\), with
\(T_n\), and (6) is the required zero-divider normalization.  For
terminal kind PE, the same conclusion holds if \(\alpha\) is already
missing on \(T_n\).  In the only remaining case, apply
\noderef{HierarchyClaimNinePENormalization} to \(P_0\), the ambient
vertex-missing hypothesis, (3), and (6).
Its first alternative proves the desired nonemptiness.  Its second
alternative supplies a new replacement presentation and a new color
which is missing on \(T_n\), on the new starred core, and at
\(y_\ell\), and for which (3) remains true.  Its
\noderef{SameGoodHierarchy} field preserves the original zero divider
count, so this output also satisfies the zero-divider disjunct below.

Thus, after replacing \(P_0,\alpha\) only in that last alternative, we
have a replacement presentation \(P\), a color still called \(\alpha\),
and all three facts
\[
\begin{split}
 &\alpha\text{ is missing on the last starred stage of }P
       \text{ and at }y_\ell,\\
 &(q_P>0\ \text{and }\alpha\in\overline{P.\mathit{coloring}}(T_{n,q_P}))
 \quad\text{or}\quad
 (q_P=0\ \text{and }\alpha\in\overline{P.\mathit{coloring}}(T_n)),\\
 &\text{the disjointness assertion (3) holds for }P.
\end{split}                                                    \tag{7}
\]

We must also transport the Situation 7.1 premise.  The
\(\mathit{sameGood}\) field of
\noderef{HierarchyReplacementPresentation}, unfolded through
\noderef{SameGoodHierarchy}, supplies
\noderef{SameHierarchyStages} between \(H_0\) and
\[
                 H=P.\mathit{counterexample}.\mathit{hierarchy}.
\]
Its fields identify the divider counts and base cores, and identify the
full stage families after transport across that count equality.  At
divider count zero the base-core equality, and at positive divider count
the equality of the last ordinary stage, show from the two branches of
\noderef{HierarchyStarredStage} that the two last starred stages are the
same exact ordered tree.  The definition
\noderef{HierarchyPathNumber} uses only \(T\), this last starred stage,
and endpoint equations for the exact edges of \(T\).  Hence its candidate
finite sets, their minima, and their empty-set fallback are identical for
\(H_0\) and \(H\).  The assumed Situation 7.1 equation therefore gives
\[
                 \noderef{HierarchyPathNumber}(H)=1.          \tag{8}
\]

For completeness, we now verify the exact augmentation structure to
which the remaining Claim 9 cases apply.  Let \(\mathcal E_P\) be the
replacement ETT in \(P\), let \(T_n\) be its last penultimate ladder
tree, and inspect its last certified
\noderef{TashkinovSeriesStep}.  Positivity of its rung again excludes
the initial constructor.

In the PE constructor, the \(\mathit{closedExtension}\) field is a
\noderef{TashkinovClosure} from \(T_n\) to the final ladder tree.
Unfolding that definition gives a
\noderef{TashkinovAugmentationSequence}; the
\(\mathit{belowFinal}\) field of
\noderef{ExtendedTashkinovTree} makes \(T\) an exact intermediate
prefix.  Since \(T_n\preccurlyeq T\) by
\(\mathit{abovePenultimate}\), applying
\noderef{TashkinovAugmentationSequenceRestrict} gives a route from
\(T_n\) to \(T\).  If the current divider count is positive, the
\(\mathit{baseStage}\) and iterated \(\mathit{nested}\) fields of
\noderef{TashkinovHierarchy} put \(T_n\preccurlyeq A\).  If the count is
zero, the PE clause of \noderef{HierarchyBase} makes \(A\) an
\noderef{OrderedTreeUnion} with \(T_n\) as its left ordered initial
segment.  Thus in either case \(T_n\preccurlyeq A\preccurlyeq T\).

In the SE and RE constructors, the closure begins at the recorded
adjoined tree \(T_n+f_n\).  Its
\noderef{TreeSequenceAdjoinEdge} certificate says that this exact tree
has length \(|T_n|+1\), begins with \(T_n\), and has new last edge the
individual connecting edge \(f_n\).  The strict
\(\mathit{abovePenultimate}\) inequality gives
\(|T_n|+1\leq |T|\).  Both \(T_n+f_n\) and \(T\) are exact ordered
initial segments of the final ladder tree: the first fact follows by
composing the one-step prefix relations in the augmentation sequence
inside \noderef{TashkinovClosure}, and the second is
\(\mathit{belowFinal}\).  Comparing their common vertex and edge entries
with \noderef{TreeSequenceInitialSegment} therefore gives
\(T_n+f_n\preccurlyeq T\), and
\noderef{TashkinovAugmentationSequenceRestrict} gives the route from this
exact seed to \(T\).  When the divider count is zero,
\noderef{HierarchyBase} identifies \(A=T_n\), so this seed supplies the
exact first tail edge.  When the count is positive, the strict hierarchy
stages give \(|T_n|+1\leq |A|\), and the same common-extension argument
puts the seed before \(A\).

Finally, by \noderef{TashkinovAugmentationSequence}, each route just
obtained is a finite reflexive-transitive chain of one-step
\noderef{TashkinovAugmentation}s.  Each such step has one additional
exact edge and retains every preceding ordered entry.  In the PE case,
and in the positive-divider SE/RE cases, induction on this finite chain
splits it at the unique member of length \(|A|\); hence every displayed
addition after \(A\) is a Tashkinov augmentation.  In the zero-divider
SE/RE case the route starts at length \(|A|+1\): its first addition is
the exact certified seed \(f_n\), and every later addition is a
Tashkinov augmentation.
To extract the endpoint statement from (8), observe first that the
empty-set fallback in \noderef{HierarchyPathNumber} would equal
\(\ell\), which is at least two by (2), and so cannot equal one.
The candidate set is therefore nonempty.  Its minimum is one, hence one
itself belongs to the filtered set; the filter condition says that every
exact edge in position at least \(|A|+1\), that is, every edge from the
second tail position onward, joins consecutive displayed vertices.
Together with the augmentation route above, this is exactly the paper's
Situation 7.1 tail.

It remains to dispatch the first use of \(\alpha\).  Consider the finite
set
\[
 \mathcal U=\{t:\operatorname{Fin}|T|:\ |A|\leq t.\mathit{val}
       \text{ and }P.\mathit{coloring}(T.\mathit{edge}(t))
                    =\alpha\}.
\]
If \(\mathcal U\) is empty, its defining negation says for every exact
tail position \(t\) that
\(P.\mathit{coloring}(T.\mathit{edge}(t))\ne\alpha\).
Together with (7), (8), and the ambient vertex-missing hypothesis, these
are precisely the hypotheses of
\noderef{HierarchyClaimNineUnusedTailColor}.  That node returns the
eligible index \(\ell-1\), so the eligible-index set is nonempty.

Otherwise let \(e_j\) be the least element of \(\mathcal U\) and define
\[
                  j=e_j.\mathit{val}-|A|+1.
\]
The lower bound in the definition of \(\mathcal U\) gives \(1\leq j\).
Also \(e_j.\mathit{val}<|T|=|A|+\ell\) by (4), so \(j\leq\ell\), and
arithmetic gives
\[
 e_j.\mathit{val}=|A|+j-1,\qquad
 P.\mathit{coloring}(T.\mathit{edge}(e_j))=\alpha.            \tag{9}
\]
Leastness says that no exact tail position strictly before \(e_j\) has
color \(\alpha\).

If \(2\leq j\), equations (7)--(9), the leastness assertion, and the
ambient vertex-missing hypothesis are exactly the premises of
\noderef{HierarchyClaimNineFirstUseAfterSeed}.  Its conclusion puts
\(j-1\) in the eligible-index set.  Otherwise \(1\leq j\) gives \(j=1\);
then (9) says that the exact first edge after \(A\) has color \(\alpha\).
Apply \noderef{HierarchyClaimNineFirstUseSeed} with (7), (8), and the
ambient vertex-missing hypothesis.  Its
conclusion is directly the nonemptiness of the eligible-index set.  The
empty and nonempty cases for \(\mathcal U\), and the two possible values
of its first subscript, exhaust all possibilities and prove Claim 9.
\end{proof}
